\chapter{数学与物理记号}
\label{chap:physics}

\section{公式、微分与自然常数}
自然常数采用正体 $\mathrm{e}$，例如
\begin{equation}
  \label{eq:exponential}
  f(x)\coloneqq \mathrm{e}^{x},\qquad
  \frac{\dd f}{\dd x}=\mathrm{e}^{x},\qquad
  \int_0^1 \mathrm{e}^{x}\dd x=\mathrm{e}-1.
\end{equation}
多行推导可以使用 \texttt{align} 环境：
\begin{align}
  (a+b)^2 &= (a+b)(a+b)\\
          &= a^2+2ab+b^2.
\end{align}
分段函数与矩阵示例：
\[
  |x|=\begin{cases}
    x, & x\geq 0,\\
    -x, & x<0,
  \end{cases}
  \qquad
  I_2=\begin{pmatrix}1&0\\0&1\end{pmatrix}.
\]
向量可写为 $\bm{r}=(x,y,z)$，偏微分可写为
$\partial f/\partial x$。

\section{Dirac 记号}
使用 \texttt{braket} 宏包提供的大写命令，括号可自动调整大小：
\[
  \Bra{\psi},\qquad
  \Ket{\psi},\qquad
  \Braket{\phi|\psi},\qquad
  \Braket{\psi|\hat{A}|\psi}.
\]
归一化态满足 $\Braket{\psi|\psi}=1$。以二能级系统为例，
\begin{equation}
  \Ket{\psi}=\frac{\Ket{0}+\Ket{1}}{\sqrt{2}},\qquad
  \hat{\rho}=\Ket{\psi}\Bra{\psi},\qquad
  \Tr(\hat{\rho})=1.
\end{equation}
当哈密顿算符 $\hat{H}$ 不显含时间时，薛定谔方程及形式解为
\begin{align}
  \mathrm{i}\hbar\frac{\partial}{\partial t}\Ket{\psi(t)}
    &=\hat{H}\Ket{\psi(t)},\\
  \Ket{\psi(t)}
    &=\mathrm{e}^{-\mathrm{i}\hat{H}t/\hbar}\Ket{\psi(0)}.
\end{align}

\section{交叉引用与表格}
使用标签引用第~\ref{chap:environments}~章或公式~\eqref{eq:exponential}，
避免手动填写编号。
\begin{table}[htbp]
  \centering
  \caption{常用物理量的记号示例}
  \label{tab:notation}
  \begin{tabular}{lll}
    \toprule
    物理量 & 记号 & SI 单位\\
    \midrule
    时间 & $t$ & $\mathrm{s}$\\
    能量 & $E$ & $\mathrm{J}$\\
    约化普朗克常数 & $\hbar$ & $\mathrm{J}\,\mathrm{s}$\\
    \bottomrule
  \end{tabular}
\end{table}

\section{链接颜色示例}

文档内部链接：参见第~\ref{chap:environments}~章。 % linkcolor
文献引用：参见文献~\cite{griffiths2018}。    % citecolor
网址链接：\url{https://www.ctan.org/}。       % urlcolor
